%!TEX root = ../thesis.tex
%*******************************************************************************
%****************************** Eight Chapter **********************************
%*******************************************************************************
\chapter{Conclusions}
\label{chap:conclusions}

% **************************** Define Graphics Path **************************
\ifpdf
    \graphicspath{{Chapter8/Figs/Raster/}{Chapter8/Figs/PDF/}{Chapter8/Figs/}}
\else
    \graphicspath{{Chapter8/Figs/Vector/}{Chapter8/Figs/}}
\fi

\epigraph{I had nothing to offer anybody except my own confusion.}{--- \textup{Jack Kerouac}}
% proofread

This chapter summarises the contributions of the research to the multiple fields of HCI, IS engineering and CI and concludes by offering short answers to the research questions introduced in the \hyperlink{questions}{Introduction}.

\section{Contributions}

\subsection{HCI}
Despite a growing body of research in HCI focusing on the design of inclusive technologies for civic participation, there is a lack of research with regard to understanding what enables citizens to become active users of data and technologies for participating in decision-making. This thesis has established an understanding of what makes data useful (i.e. effective use) for non-expert citizens and communities, contributing new knowledge about the relationships between data, communities, civic participation, advocacy and action. Through PADRE, which is a common approach for studies in HCI, and applying engineering sprints to develop and deploy technologies, this study actively responded to community problem-solving activities. Furthermore, by analysing the way people made use of community data sources, new understandings of civic advocacy and action were uncovered, which were condensed into a CAF that helped to discuss advocacy efforts and plan for future activities relating to local action by citizens (Chapter \ref{chap:sensemystreet-eval}). However, the analysis and evaluation of these technologies also revealed the temporality of such solutions responding to ever-changing communities and their activities. Building on community commissioning research \citep{garbett_designing_2017}, this thesis extended the digital civics approach \citep{olivier_digital_2015,vlachokyriakos_digital_2016,asad_tap_2017} with the inclusion of designing for building community capacity as a step towards the successful transfer of ownership and the adoption of community information resources (Chapters \ref{chap:sensemystreet-eval} and \ref{chap:cdi}), in addition to presenting a deeper understanding of the role of expert professionals and instructional designers or researchers in all of this.

\subsection{IS}
Taking a human-centred design approach to designing and developing community IS, this research documented the design processes and also provided two open-source platforms for IS: (1) the SMS toolkit (Chapter \ref{chap:sensemystreet}), which is a set of tools and platforms that enable communities to use scientific-grade environmental monitors to produce and commission data about their neighbourhoods and (2) Data:In Place (Chapter \ref{chap:datainplace}), an open-source web platform built to support citizens in accessing, interpreting and using open data for the purposes of civic advocacy and action. Furthermore, through the community problem-solving activities and the user-centred design processes, the thesis contributed new ways to accommodate more social agendas into the design and development of future IS.

\subsection{CI}
This thesis provided a new model of CDI for the field of CI that helps discuss and design community resources for knowledge making and civic action through the effective use of data by citizens (Chapter \ref{chap:cdi}). CDI extends the view relating to the design of CI tools from focusing only on the access and instrumentation of control \citep{stillman_community_2009} to looking at how technologies could be used also to communicate ideas, share knowledge, create community links and increase capacities for the effective use of data by communities. The thesis also contributed findings from the practical use of CDI as a technical solution and social facilitation method that support communities making effective use of data for civic advocacy and action (Chapter \ref{chap:cdi_in_action}). This provides practical examples for the key elements of CDI to take, enabling researchers to better adopt the model for further research in CI and CDI.

\section{Answering the Research Questions}
\subsubsection{What is involved in citizens creating actionable knowledge using data?}
The assumption about a lack of skills, access or tools to use data being the only reason for communities not leveraging that data is not completely true. However, what was learned from the research was that building better systems for the access and use of data does not necessarily increase data's actionability by communities. In the same way, providing new ways for citizens to actively produce or commission data will not guarantee the effective use of that data by citizens. The usefulness of data for citizens is dictated by the activities they plan to carry out and what they want to achieve. Citizens need a purpose for using the data, and only then can it become useful to them and be put in action. Furthermore, such purposes need to be integrated into the investigation and issue enquiry process at an early stage. 

\subsubsection{How do we build and configure tools and processes to support knowledge making and enable the effective use of data by citizens in civic advocacy and action?}
In order to enable communities to effectively use data for the local benefit, we need to move away from building tools solely for accessing and using data and start looking at ways to build tools for social facilitation around the use of data. This includes providing training, methods and digital tools to help people articulate their concerns, communicate and discuss ideas, understand what resources are needed and plan steps for action. Furthermore, some of these tools are already used by communities for other purposes, which could be reconfigured to serve communities' problem-solving agendas. Indeed, there is a need for increased capacities for communities and either the transfer of or support of data science skills. However, this does not mean coming in and doing the work for the communities. In order to maintain the engagement and use of community information resources and solutions, the community needs to be able to take ownership and, through increasing community links, adapt to a constantly changing social economy.